% !PRINTED 611
\weberpage{611}
Wir bilden jetzt den Körper:

\weq{16}{\mathfrak{K}'' = \Omega(\xi, \xi') = \Omega(\xi''),}

indem wir

\weq{17}{\xi'' = \alpha\xi + \alpha'\xi'}

setzen, und bezeichnen ein Primideal ersten Grades dieses Körpers,
das in einem $\mathfrak{p}_1'$ aufgeht, mit $\mathfrak{P}''$.

$\mathfrak{P}, \mathfrak{P}'$ seien, wie oben, die durch $\mathfrak{P}''$ teilbaren Primideale in
$\mathfrak{K}, \mathfrak{K}'$. Dann ist nach (15) und (17)
\[
\xi''^{p} \equiv \xi'' \quad (\mathrm{mod}\ \mathfrak{P}''),
\]
und folglich ist $\mathfrak{K}''$ ein Klassenkörper von $A'$.

Demnach ist $\mathfrak{K}''$ mit $\mathfrak{K}'$ identisch; $\xi'$ ist eine primitive Zahl
von $\mathfrak{K}''$, und folglich kann jede Zahl in $\mathfrak{K}''$, also auch $\xi$, rational
durch $\xi'$ ausgedrückt werden.

Daraus folgt:

11. \gesperrt{Enthält eine Zahlgruppe $A$ eine andere Zahlgruppe
 $A'$ als Teiler von endlichem Index, und existieren
 die Klassenkörper $\mathfrak{K}(A)$, $\mathfrak{K}(A')$, so ist $\mathfrak{K}(A)$
in $\mathfrak{K}(A')$ enthalten.}

\begin{center}
\textbf{§ 165. Primideale in den Klassen.}
\end{center}

Wir kehren jetzt zu den Funktionen $Q$ zurück, die wir in
§ 163, (16) durch unendliche Reihen und in § 163, (18) durch
unendliche Produkte:

\weq{1}{Q_{\varkappa} = \Pi\, \frac{1}{1 - \chi_{\varkappa}(\mathfrak{p})\,N(\mathfrak{p})^{-s}}}

dargestellt haben, worin $\chi_{\varkappa}$ einen der Charaktere der Gruppe
bedeutet und $\mathfrak{p}$ alle Primideale in $\overline{O}$ durchläuft. Den Logarithmus
 dieses Produktes entwickeln wir in folgender Weise:

\weq{2}{\log Q_{\varkappa} = \sum \frac{\chi_{\varkappa}(\mathfrak{p})}{N(\mathfrak{p})^{s}} + \frac{1}{2}\sum \frac{\chi_{\varkappa}(\mathfrak{p}^2)}{N(\mathfrak{p})^{2s}} + \frac{1}{3}\sum \frac{\chi_{\varkappa}(\mathfrak{p}^3)}{N(\mathfrak{p})^{3s}} + \cdots}

Das Vielfache von $2\pi i$, das in dem $\log Q$ nicht näher definiert
ist, brauchen wir nicht zu kennen, da es sich mit $s$ nur unstetig
ändern könnte, und da rechts und links stetige Funktionen von
$s$ stehen, solange $s > 1$ ist, so ändert sich dieses Vielfache überhaupt
 nicht mit $s$.

Es sei $\overline{A}_i$ eine der Klassen der Idealgruppe $\overline{A}$ und $\overline{A}_i'$ die
entgegengesetzte Klasse. Dann ist nach (12), (13), (14), § 163:
\[
\overset{k}{\sum} \chi_{\varkappa}(\overline{A}_i')\,\chi_{\varkappa}(\mathfrak{a}) = h \ \text{oder} = 0,
\]
%CONT>>
